\Needspace{13\baselineskip}
\Article{18}

The method of the preceding article supplied the absolute values of the functions $T$, $U$, and specified the conditions under which they must be given equal or opposite signs; but it does not yet determine the signs themselves.\quad For the case where $k=1$, we shall supply what is missing as follows.

Let us put $\rho=\cos\frac12\omega+i\sin\frac12\omega$, so that $r=\rho\rho$. Since $\rho^n=-1$, it is clear that equation $[8]$ can be written as

\[
\begin{aligned}
W&=2(1+\rho^{n-2})(1+\rho^{-4})(1+\rho^{n-6})(1+\rho^{-8})\\&\qquad\ldots(1+\rho^{-n+4})(1+\rho^2)
\end{aligned}
\]

\Needspace{6\baselineskip}
or, with the factors arranged in a different order,

\[
\begin{aligned}
W&=2(1+\rho^2)(1+\rho^{-4})(1+\rho^6)(1+\rho^{-8})\\&\qquad\ldots(1+\rho^{-n+4})(1+\rho^{n-2})
\end{aligned}
\]

Now we have

\[
\begin{aligned}
1+\rho^2&=2\rho\cos\frac12\omega\\
1+\rho^{-4}&=2\rho^{-2}\cos\omega\\
1+\rho^{-6}&=2\rho^3\cos\frac32\omega\\
1+\rho^{-8}&=2\rho^{-4}\cos2\omega,\;\text{etc.}
\end{aligned}
\]

up to

\[
\begin{aligned}
1+\rho^{-n+4}&=2\rho^{-\frac12n+2}\mathord{,}\cos(\frac14n-1)\omega\\
1+\rho^{n-2}&=2\rho^{\frac12n-1}\cos(\frac14n-\frac12)\omega
\end{aligned}
\]

Thus we have

\SourceLeaf{41}

\[
\begin{aligned}
W&=2^{\frac12n}\rho^{\frac14n}\cos\frac12\omega\cos\omega\,\frac32\omega\\&\qquad\dotsfour\cos(\frac14n-\frac12)\omega
\end{aligned}
\]

\Needspace{8\baselineskip}
The cosines entering this product are evidently all positive, while the factor $\rho^{\frac14n}$ is $=\cos45^0+i\sin45^0=(1+i)\surd\frac12$.\quad From this we conclude that $W$ is the product of $1+i$ and a positive real quantity, so that necessarily

\[
\begin{aligned}
W=(1+i)\surd n,\quad T=+\surd n,\quad U=+\surd n
\end{aligned}
\]

\Needspace{20\baselineskip}
\Article{19}

It will be worthwhile to collect all the summations developed so far into a single overview here.\quad In general, we have

\[
\begin{array}{c|c|c}
T=&U=&\text{for }n\text{ of the form}\\\hline
\pm\surd n&\pm\surd n&4\mu\\
\pm\surd n&0&4\mu+1\\
0&0&4\mu+2\\
0&\pm\surd n&4\mu+3
\end{array}
\]

and, when we set $k$ $=1$, the radical must be given the positive sign.\quad What we observed by induction in Article 3 for prime values of $n$ has therefore now been proved with full rigor. Nothing remains except to show how to determine the signs in every case, for any values of $k$.\quad But before we can tackle this task in full generality, we must first consider more closely the cases where $n$ is a prime number or a power of a prime number.

\Needspace{14\baselineskip}
\Article{20}

First let $n$ be an odd prime number. From what we set out in Article 10, it is clear that $W=1+2\Sigma r^a=1+2\Sigma R^{ak}$ if we put $R=\cos\omega+i\sin\omega$, with $a$, as there, denoting indefinitely all the quadratic residues of $n$ lying between $1$ and $n-1$.\quad If we likewise use $b$ to denote indefinitely all the quadratic nonresidues between the same limits, we readily see that all the numbers $ak$ are congruent modulo $n$ either to all the $a$ or to all the $b$ (without regard to order), according as $k$ is a residue or a nonresidue.\quad Thus, in the first case, we have

\SourceLeaf{42}

\[
\begin{aligned}
W=1+2\Sigma R^a=1+R+R^4+R^9+\text{etc.}+R^{(n-1)^2}
\end{aligned}
\]

and hence $W=+\surd n$ if $n$ is of the form $4\mu+1$, and $W=+i\surd n$ if $n$ is of the form $4\mu+3$.

In the other case, where $k$ is a nonresidue of $n$, we have

\[
\begin{aligned}
W=1+2\Sigma R^b
\end{aligned}
\]

\Needspace{12\baselineskip}
Since all the $a$, $b$ evidently make up the complete collection of numbers $1$, $2$, $3$ $\ldots$, and therefore

\[
\begin{aligned}
\Sigma R^a+\Sigma R^b=R+R^2+R^3+\text{etc.}+R^{n-1}=-1
\end{aligned}
\]

we obtain

\[
\begin{aligned}
W=-1-2\Sigma R^a=-(1+R+R^4+R^9+\text{etc.}+R^{(n-1)^2})
\end{aligned}
\]

and hence $W=-\surd n$ if $n$ is of the form $4\mu+1$, and $W=-i\surd n$ if $n$ is of the form $4\mu+3$.

\Needspace{20\baselineskip}
From this we conclude:

\emph{first}, if $n$ is of the form $4\mu+1$ and $k$ is a quadratic residue of $n$,

\[
\begin{aligned}
T=+\surd n,\quad U=0
\end{aligned}
\]

\emph{second}, if $n$ is of the form $4\mu+1$ and $k$ is a nonresidue of $n$,

\[
\begin{aligned}
T=-\surd n,\quad U=0
\end{aligned}
\]

\emph{third}, if $n$ is of the form $4\mu+3$ and $k$ is a residue of $n$,

\[
\begin{aligned}
T=0,\quad U=+\surd n
\end{aligned}
\]

\emph{fourth}, if $n$ is of the form $4\mu+3$ and $k$ is a nonresidue of $n$,

\[
\begin{aligned}
T=0,\quad U=-\surd n
\end{aligned}
\]
